\input{../tex-inputs/latex-leseansicht-vorspann}

\section[Arthur und Olga Schnitzler an Gustav Schwarzkopf, 30. 5. 1910]{L04366 Arthur und Olga Schnitzler an Gustav Schwarzkopf, 30. 5. 1910}
\nopagebreak\mylabel{L04366v}
\rehead{ }\normalsize\beginnumbering\briefempfaengerindex{Schwarzkopf, Gustav@\textsc{Schwarzkopf, Gustav}!zzzSchnitzler, Olga@\emph{von Olga Schnitzler}!1910-05-302@{30. 5. 1910}|(be}\briefempfaengerindex{Schwarzkopf, Gustav@\textsc{Schwarzkopf, Gustav}!zzzSchnitzler, Arthur@\emph{von Arthur Schnitzler}!1910-05-302@{30. 5. 1910}|(be}
\toendnotes[C]{\smallbreak\pagebreak[2]}
\correspDesc{Versand  durch Arthur Schnitzler, Olga Schnitzler am 30. 5. 1910 in Genf
\newline{}Übermittlung  am 30. 5. 1910 in Genf
\newline{}Erhalt  durch Gustav Schwarzkopf im Zeitraum [31. 5. 1910 – 4. 6. 1910?] in Wien}\toendnotes[C]{\smallbreak}
\Standort{DLA, A:Schnitzler, HS.1985.1.1897.}
\physDesc{Bildpostkarte, 111 Zeichen
\newline{}Handschrift Arthur Schnitzler: Bleistift, deutsche Kurrent
\newline{}Handschrift Olga Schnitzler: Bleistift, lateinische Kurrent
\newline{}Versand: Stempel: »\nobreak{}\oindex{Genf@\textbf{Genf}|pwk}Genève Exp. Lettr{[}es{]}, 30. V. 10, 12\nobreak{}«.  }\pstart{}{\pb}{[}hs. O. Schnitzler:{]} Autriche\oindex{Österreich@\textbf{Österreich}|pw}\pend{}\pstart{}{[}hs. A. Schnitzler:{]} Hrn \textsc{Gustav Schwarzkopf}\pend{}\pstart{}Wien I\oindex{I., Innere Stadt@\textbf{I., Innere Stadt}|pw}\pend{}\pstart{}\textsc{Tiefer Graben 23}\oindex{Wien@\textbf{Wien}!I., Innere Stadt@\textbf{I., Innere Stadt}!Tiefer Graben 23@\textbf{Tiefer Graben 23}|pw}\pend{}{\bigskip}
\pstart
           \noindent{}\centering{}{\pb}\textcolor{gray}{\textbf{Genève\oindex{Genf@\textbf{Genf}|pw}. – Pont du Mont-Blanc\oindex{Mont Blanc Brücke@\textbf{Mont Blanc Brücke}|pw} et Quai des Bergues\oindex{Quai des Bergues@\textbf{Quai des Bergues}|pw}.}}\pend
           \vspace{1em}
\pstart
           \raggedleft{}{\pb}\textsc{Genf\oindex{Genf@\textbf{Genf}|pw}}, 30. 5. 10\pend
           \vspace{0.5em}
\pstart
           Herzliche Grüße,{\\[\baselineskip]}Ihr \spacefill\mbox{Arthur}\pend
           \leftskip=0em{}\selectlanguage{ngerman}\vspace{1em}
\pstart
           \noindent{}{[}hs. O. Schnitzler:{]} Auf Wiedersehen!\pend
           \pstart \spacefill\mbox{Olga.}\pend{}\selectlanguage{ngerman}\endnumbering\briefempfaengerindex{Schwarzkopf, Gustav@\textsc{Schwarzkopf, Gustav}!zzzSchnitzler, Olga@\emph{von Olga Schnitzler}!1910-05-302@{30. 5. 1910}|)be}\briefempfaengerindex{Schwarzkopf, Gustav@\textsc{Schwarzkopf, Gustav}!zzzSchnitzler, Arthur@\emph{von Arthur Schnitzler}!1910-05-302@{30. 5. 1910}|)be}\mylabel{L04366h}
\begin{anhang}
\end{anhang}\newcommand{\dateiname}{L04366}\newcommand{\titel}{Arthur und Olga Schnitzler an Gustav Schwarzkopf, 30. 5. 1910}\newcommand{\editorInnen}{Herausgegeben von Jahnke, SelmaMüller, Martin Anton}\input{../tex-inputs/latex-leseansicht-abspann}
